\documentclass[aspectratio=169,10pt]{beamer}
\usetheme[numbering=fraction,progressbar=frametitle]{metropolis}
\usepackage{fontspec}
\setsansfont{Latin Modern Sans}
\usepackage{amsmath,amssymb,mathtools,graphicx,booktabs}
\usepackage{polyglossia}
\setdefaultlanguage{french}
\setbeamertemplate{navigation symbols}{}
\setbeamerfont{frametitle}{size=\large}
\setbeamerfont{normal text}{size=\normalsize}
\setlength{\parskip}{.45em}
\hypersetup{pdfauthor={},pdfsubject={Lecture slides},pdfkeywords={},pdftitle={Dérivation et tangentes}}
\providecommand{\N}{\mathbb{N}}
\providecommand{\Z}{\mathbb{Z}}
\providecommand{\D}{\mathbb{D}}
\providecommand{\Q}{\mathbb{Q}}
\providecommand{\R}{\mathbb{R}}
\providecommand{\C}{\mathbb{C}}
\providecommand{\vect}[1]{\overrightarrow{#1}}
\providecommand{\vi}{\vec{\imath}}
\providecommand{\vj}{\vec{\jmath}}
\providecommand{\vk}{\vec{k}}
\providecommand{\norm}[1]{\left\lVert #1\right\rVert}
\providecommand{\vc}[2]{\begin{pmatrix}#1\\#2\end{pmatrix}}
\providecommand{\deter}[2]{\det\left(#1,#2\right)}
\providecommand{\abs}[1]{\left\lvert #1\right\rvert}
\providecommand{\dd}{\,\mathrm{d}}
\providecommand{\e}{\mathrm{e}}
\providecommand{\Card}{\operatorname{Card}}
\providecommand{\E}{\operatorname{E}}
\providecommand{\V}{\operatorname{V}}
\providecommand{\degre}{^{\circ}}
\providecommand{\un}[1]{\,\mathrm{#1}}
\providecommand{\sci}[2]{#1\times10^{#2}}
\providecommand{\pgcd}{\operatorname{PGCD}}
\providecommand{\ppcm}{\operatorname{PPCM}}
\title{Dérivation et tangentes}
\subtitle{Cours et exemples}
\date{}
\author{}
\begin{document}
\begin{frame}[plain]
\vfill
{\usebeamerfont{title}\inserttitle\par}
\medskip
{\usebeamerfont{subtitle}\insertsubtitle\par}
\medskip
{\usebeamercolor[fg]{progress bar}\rule{\textwidth}{.5pt}}
\vfill
\end{frame}
\begin{frame}{Objectifs}
Calculer un nombre dérivé, dériver les fonctions usuelles et utiliser le signe de la dérivée pour déterminer les variations.
\end{frame}
\begin{frame}{Du taux de variation au nombre dérivé}
Si \(a\) est intérieur au domaine et si \(\lim_{h\to0}\dfrac{f(a+h)-f(a)}h\) existe et est un réel, \(f\) est dérivable en \(a\) et cette limite est \(f'(a)\). La fonction dérivée associe à chaque point où ce nombre existe sa valeur \(f'(a)\).
\end{frame}
\begin{frame}{Du taux de variation au nombre dérivé}
Si \(f\) est dérivable en \(a\), la tangente au point \((a;f(a))\) a pour équation \(y=f'(a)(x-a)+f(a)\).
\end{frame}
\begin{frame}{Du taux de variation au nombre dérivé}
La droite passe par \((a;f(a))\) et a pour coefficient directeur la limite des pentes des sécantes, \(f'(a)\). La forme point-pente donne l’équation.
\end{frame}
\begin{frame}{Du taux de variation au nombre dérivé}
La dérivabilité en un point entraîne la continuité en ce point. La réciproque est fausse.
\end{frame}
\begin{frame}{Du taux de variation au nombre dérivé}
Pour \(h\ne0\), \(f(a+h)-f(a)=h\dfrac{f(a+h)-f(a)}h\). Le second facteur tend vers un réel, donc le produit tend vers \(0\). Ainsi \(f(a+h)\to f(a)\). Pour \(|x|\) en \(0\), les taux sont \(-1\) à gauche et \(1\) à droite : continuité, mais pas dérivabilité.
\end{frame}
\begin{frame}{Dérivées usuelles et opérations}
Pour \(u,v\) dérivables et \(\lambda\in\mathbb{R}\),
\[
\begin{aligned}
(u+v)'&=u'+v',& (\lambda u)'&=\lambda u',\\
(uv)'&=u'v+uv',& \left(\dfrac uv\right)'&=\dfrac{u'v-uv'}{v^2}\quad(v\ne0).
\end{aligned}
\]
Si \(u\) est dérivable et \(F\) dérivable aux valeurs de \(u\), \((F\circ u)'=(F'\circ u)u'\). Les formules générales sont admises, sauf la justification du produit ci-dessous.
\end{frame}
\begin{frame}{Dérivées usuelles et opérations}
Ajouter et retrancher \(u(a+h)v(a)\) dans la différence du produit donne
\[
\dfrac{u(a+h)v(a+h)-u(a)v(a)}h
=u(a+h)\dfrac{v(a+h)-v(a)}h
+v(a)\dfrac{u(a+h)-u(a)}h.
\]
La continuité de \(u\) et les limites des taux donnent \(u(a)v'(a)+v(a)u'(a)\).
\end{frame}
\begin{frame}{Dérivées usuelles et opérations}
Pour \(f(x)=\dfrac{x^2+1}{x-1}\), \(x\ne1\),
\[
f'(x)=\dfrac{2x(x-1)-(x^2+1)}{(x-1)^2}
=\dfrac{x^2-2x-1}{(x-1)^2}.
\]
Pour \(g(x)=\sqrt{3x+2}\), \(g'(x)=3/(2\sqrt{3x+2})\) sur \(]-2/3;+\infty[\).
\end{frame}
\begin{frame}{Variations et extrema}
Sur un intervalle, une fonction dérivable dont la dérivée est positive ou nulle est croissante; si elle est négative ou nulle, elle est décroissante. Une dérivée strictement positive donne une croissance stricte. Si elle est nulle partout sur l’intervalle, la fonction est constante.
\end{frame}
\begin{frame}{Variations et extrema}
Un zéro isolé de \(f'\) n’impose pas un extremum : \(x^3\) a une dérivée nulle en \(0\) et reste croissante. Un changement de signe de \(+\) vers \(-\) donne un maximum local; de \(-\) vers \(+\), un minimum local. Un extremum au bord d’un intervalle ne nécessite pas une dérivée nulle.
\end{frame}
\begin{frame}{Variations et extrema}
Déterminer le domaine; calculer et factoriser la dérivée; étudier son signe sur chaque intervalle; calculer les valeurs aux points critiques et les limites aux bornes; tracer les flèches de variation.
\end{frame}
\begin{frame}{Variations et extrema}
Repérer la fonction extérieure et l’expression intérieure. Pour \((u(x))^n\), dériver la puissance puis multiplier par \(u'(x)\); pour \(\sqrt{u(x)}\), exiger \(u(x)>0\) avant d’utiliser la formule; pour un quotient, écrire le numérateur de la dérivée avant de le développer. Toujours annoncer le domaine de dérivation.
\end{frame}
\begin{frame}{Variations et extrema}
Pour \(F(x)=\sqrt{2x+3}\), le domaine de définition est \([-3/2;+\infty[\). La formule de dérivation donne, pour \(x>-3/2\),
\[
F'(x)=\dfrac{2}{2\sqrt{2x+3}}=\dfrac1{\sqrt{2x+3}}.
\]
L’extrémité est dans le domaine de définition mais pas dans ce domaine de dérivation. Pour \(G(x)=(x^2+1)^3\), les deux fonctions sont dérivables sur \(\mathbb{R}\) et \(G'(x)=3(x^2+1)^2\times2x\).
\end{frame}
\begin{frame}{Variations et extrema}
Après avoir calculé \(f'\), déterminer son signe sur le domaine physique du problème. Un changement de \(+\) à \(-\) prouve un maximum local; pour un maximum global, comparer tous les intervalles et les éventuelles bornes incluses. Donner la valeur maximale, l’endroit où elle est atteinte et les unités. Un point où la dérivée est nulle est seulement un candidat.
\end{frame}
\begin{frame}{Variations et extrema}
Pour un rectangle de périmètre \(20\) m, si \(x\) est un côté en mètres, l’autre vaut \(10-x\) et \(0<x<10\). L’aire \(A(x)=10x-x^2\) a une dérivée \(10-2x\), qui change de signe en \(5\). Le maximum vaut \(25\ \mathrm{m}^2\), pour un carré. La forme \(A(x)=25-(x-5)^2\) fournit une seconde preuve globale et précise l’égalité.
\end{frame}
\begin{frame}{Variations et extrema}
Pour rédiger sa recherche, calculer séparément les coordonnées du point et la pente, puis utiliser la forme point-pente. Comparer la courbe à cette droite signifie étudier le signe de la différence des ordonnées; regarder seulement le dessin ne prouve pas ce signe.
\end{frame}
\end{document}
